\documentclass[11pt]{article}

\usepackage[utf8]{inputenc}
\usepackage[T1]{fontenc}
\usepackage[margin=1in]{geometry}
\usepackage{booktabs}
\usepackage{hyperref}
\usepackage{url}

\hypersetup{colorlinks=true, linkcolor=blue, urlcolor=blue, citecolor=blue}

\title{CycloneNet: A Reproducible Physics-Guided Pipeline for Forensic
Rapid-Intensification Analysis --- and an Honest Audit of What It Can and
Cannot Do}
\author{Estefano Senhor Ferreira\\
\small Software Engineer, Independent Researcher\\
\small \href{mailto:estefano.senhor@gmail.com}{estefano.senhor@gmail.com}
\,\textperiodcentered\,
\url{https://github.com/estefano-ferreira/cyclone-net}}
\date{}

\begin{document}

\maketitle

\begin{center}
\small\itshape Revised, corrected manuscript (supersedes ``CycloneNet V2:
\ldots{} Atmospheric Singularity Mapping'', Feb 2026; this revision, Jul 2026,
adds the full 1980--2023 data release and results). See \S8 for the
correction record.
\end{center}

\begin{abstract}
CycloneNet is an open-source, configuration-driven pipeline for forensic
(hindcast) analysis of tropical-cyclone rapid intensification (RI) from ERA5
reanalysis. It couples a 3D-CNN RI classifier with an interpretable spatial
output (``FuelMap'') and physics-guided weak-constraint losses.

This paper is deliberately \textbf{not} a claim of novelty. The individual
components (3D-CNNs for RI, physics-informed/interpretable ML for tropical
cyclones, ocean-heat predictors of RI) are established in the literature, and
the central physical question --- \emph{where is the ocean energy that fuels a
storm?} --- is already addressed operationally by Tropical Cyclone Heat
Potential (TCHP)/ocean-heat-content products and by coupled ocean--atmosphere
models (HWRF, HAFS) that compute air--sea enthalpy flux directly. CycloneNet
does not advance the state of the art in RI skill, nor does it discover a new
physical mechanism.

What this paper contributes is twofold and modest: (1) a \textbf{transparent,
auditable, reproducible engineering artifact} --- a hybrid physics+ML hindcast
pipeline with leakage-safe splits, train-only normalization, unit tests, and a
counterfactual causal test; and (2) an \textbf{honest audit} that yields
cautionary, mostly \textbf{negative} results:
\begin{itemize}
  \item The learned FuelMap does \textbf{not} localize the ocean energy proxy
        (TCHP) better than a trivial ``predict the storm centre'' baseline.
  \item Adding a surface-altimetry ocean channel (absolute dynamic topography,
        ADT) --- itself a known TCHP proxy --- does \textbf{not} measurably
        improve RI skill in this small-data regime.
  \item The ``equation-consistency'' physics loss is \textbf{near-degenerate}
        by construction.
  \item A prior release of this pipeline reported the model as
        ``physics-guided'' while the physics losses were, in fact, inactive
        (weight 0) in the released configuration.
\end{itemize}

The audit's central reproducibility failure --- headline metrics (ROC-AUC 0.83)
obtained on an unreleased dataset --- is resolved in this revision by releasing
the full \textbf{1980--2023} dataset (16{,}780 events, 802 RI positives, 992
storms) built by a windowed, checksummed pipeline, and retraining
deterministically on it. On the held-out test split (2{,}679 events, 115 RI
positives) the classifier achieves \textbf{ROC-AUC 0.796 [95\% CI
0.753--0.837]} and \textbf{PR-AUC 0.251 [0.179--0.331]} at the test split's
4.3\% prevalence --- the first version of this project whose skill is statistically
distinguishable from chance, and the confirmation-by-intervention of the
audit's diagnostic that sample size, not architecture, was the bottleneck.
The negative results on spatial attribution stand and are strengthened by two
additional controls. We present this as a reproducible baseline and as a
caution against overclaiming in physics-guided cyclone ML.
\end{abstract}

\section{Introduction}

Tropical cyclones are heat engines: they convert heat stored in the upper
ocean into kinetic energy through air--sea enthalpy flux, with the rate of
extraction governed by the wind field and modulated by subsurface ocean heat
content (Emanuel's potential-intensity theory). Predicting rapid
intensification (RI; a $\geq$30\,kt increase in 24\,h) is an active, mature
operational and research problem.

CycloneNet asks a \emph{forensic} variant of this question: for a historical
storm episode, which environmental conditions in the reanalysis are most
consistent with intensification-supportive thermodynamics? This is a
diagnostic/interpretive framing, not a forecasting one. We stress at the
outset that this framing, and the methods used to pursue it, are \textbf{not
new} (\S2). The purpose of this paper is to document the system honestly ---
including what an audit shows it cannot do --- rather than to assert a
contribution it does not make.

\section{Positioning relative to existing work (what already exists)}

It is important to state plainly what CycloneNet does \emph{not} originate:

\begin{itemize}
  \item \textbf{Deep learning for RI} is established, including convolutional
        models over the North Atlantic and East Pacific and
        ensemble/transformer approaches that match or exceed the skill
        reported here.
  \item \textbf{Ocean heat content / altimetry for RI} is foundational
        operational oceanography: TCHP can be estimated from
        altimetry-derived sea-surface-height fields and is already a
        high-performing predictor in the operational SHIPS Rapid
        Intensification Index. The ADT $\leftrightarrow$ TCHP relationship
        CycloneNet ``rediscovers'' (\S5) is decades old.
  \item \textbf{Physics-informed and interpretable ML for TC intensity}
        (spatial attention, physics-informed networks, interpretable
        transformers) is an active subfield.
  \item \textbf{Identifying and quantifying the ocean energy source} of a
        storm is operational: TCHP/OHC products (NOAA/AOML, NESDIS, Navy
        NCODA) map the available fuel, and coupled models (HWRF since 2007;
        HAFS, official 2023) compute the air--sea enthalpy-flux field and the
        storm-induced cold wake every cycle.
\end{itemize}

CycloneNet occupies none of this as new ground. Its only non-trivial
differentiator is engineering discipline (auditability, reproducibility,
testing), which is a matter of execution quality, not scientific contribution.

\section{Data}

\begin{itemize}
  \item \textbf{ERA5} (Copernicus CDS), 0.25$^{\circ}$: SST, MSLP, 10\,m
        winds; optionally 2\,m temperature and dewpoint for bulk heat-flux
        diagnostics. The \textbf{public release covers 1980--2023}, North
        Atlantic sector season months. Because the raw monthly files
        ($\sim$60\,GB) exceed local storage, the dataset is built by a
        \textbf{windowed pipeline} (download $\rightarrow$ extract
        $\rightarrow$ verify $\rightarrow$ discard raw), with per-window
        provenance manifests recording SHA-256 checksums, per-event
        extraction outcomes, and the deletion record --- every released cube
        is auditable back to its source file.
  \item \textbf{IBTrACS} best-track for storm centres, intensities, and RI
        labels ($\Delta v \geq 30$\,kt/24\,h). Event identifiers embed the
        storm ID (\texttt{era5\_\{timestamp\}\_\{SID\}}), so concurrent storms
        at the same timestamp never collide.
  \item \textbf{TCHP} (NOAA/AOML ERDDAP) and \textbf{sea-level anomaly / ADT}
        (Copernicus DUACS) for external ocean validation. Note: the AOML
        gridded TCHP product covers 2022--present only; pre-2022 gridded TCHP
        is not publicly available, which limits validation to 2022--2023.
\end{itemize}

For each event, a $40\times40$ ($\approx$10$^{\circ}\times$10$^{\circ}$)
spatio-temporal cube over five 6-hourly steps ($t_0 \ldots t_{-24\mathrm{h}}$)
is extracted with nine input channels (SST, MSLP, U10, V10, wind speed,
vorticity, divergence, $|\nabla \mathrm{MSLP}|$, SST anomaly). Splits are
storm-level and \textbf{hash-deterministic}: each storm is assigned to
train/val/test by the SHA-256 hash of its identifier, so adding new storms
never reassigns existing ones (the previous seeded-shuffle scheme reassigned
136 of 155 test events when the dataset grew --- a leakage mechanism
eliminated in this revision). Normalization statistics are computed on the
training split only. The full dataset yields \textbf{16{,}780 valid events
from 992 storms (802 RI positives, 4.8\% global prevalence)}; train/val/test =
11{,}150/2{,}951/2{,}679 events with 545/142/115 positives (per-split
prevalence 4.3--4.9\%).

\section{Methods}

\textbf{Model.} A two-layer 3D-CNN produces (i) an RI logit and auxiliary
12/24\,h intensity-change regressions via global pooling, and (ii) a 2D
FuelMap logit map via a $1\times1$ head, from which a soft-argmax yields
continuous predicted coordinates.

\textbf{Physics-guided losses (now correctly active).} Controlled by explicit
weights in \texttt{config.yaml}:
\begin{itemize}
  \item \emph{FuelMap--prior alignment} (KL): aligns the FuelMap distribution
        with a heuristic physical prior
        $P = \mathrm{ReLU}(\mathrm{SST}_{\mathrm{anom}})\cdot
        \mathrm{wind}\cdot(1+\mathrm{ReLU}(-\mathrm{div}))$. The prior is
        explicitly heuristic.
  \item \emph{Forward constraint}: an energy score from the overlap of
        FuelMap and prior is mapped to the 24\,h intensity change,
        supervising ``localized surface energy $\rightarrow$
        intensification''.
  \item \emph{TV / L1} regularizers on the FuelMap.
  \item \emph{Equation consistency} (off by default): we document that this
        term is \textbf{near-degenerate} --- it compares
        vorticity/divergence recomputed from the wind field against stored
        channels derived from the same wind field; empirically, enabling it
        changes the loss by $<$0.001\%.
\end{itemize}

\textbf{Optional ADT ocean channel.} ADT can be sampled onto each event grid
and appended as an extra input channel (with a missingness mask where ocean
coverage is absent).

\textbf{Threshold.} Selected on validation by a configurable policy; the
default targets the project's high-recall forensic regime
(highest-precision threshold meeting a recall target).

\textbf{Causal ablation (with control).} To test whether the model's RI
prediction depends on the region the FuelMap identifies, we occlude the
top-$k$\% FuelMap pixels and compare the drop in predicted RI probability
against occluding an equal-size \emph{low-fuel control} region, using a
paired test.

\textbf{Spatial validation (with baseline).} Where TCHP is available, we
compare the distance from the FuelMap peak to the TCHP peak against the
distance from the \textbf{storm centre} to the TCHP peak --- the model only
demonstrates spatial skill if it beats this baseline.

\section{Results}

\subsection{Classification (reproducible from public data)}

On the held-out 1980--2023 test split ($n = 2{,}679$; \textbf{115 RI
positives}), the physics-guided model achieves \textbf{ROC-AUC 0.796 [95\% CI
0.753--0.837]} and \textbf{PR-AUC 0.251 [0.179--0.331]} against the test
split's 4.3\% prevalence (5.8$\times$ chance), with recall 0.852 and precision 0.070 at the
validation-selected high-recall operating point (Brier 0.037, ECE 0.011;
bootstrap CIs, 10{,}000 resamples). Both confidence intervals sit entirely
above chance. Training is deterministic --- the released checkpoint was
reproduced digit-for-digit from scratch --- and every number in this section
is computable from the public repository and dataset, resolving the
reproducibility failure recorded in \S8.

The trajectory across dataset generations is itself a result. The same
architecture, losses, and hyperparameters produce:

\begin{center}
\begin{tabular}{lccc}
\toprule
Dataset & Test positives & ROC-AUC [95\% CI] & PR-AUC $\div$ prevalence \\
\midrule
2020--2023 (original release) & 9 & 0.590 [0.334--0.842] & 2.1$\times$ \\
2020--2023, collision-fixed IDs & 13 & 0.614 [0.453--0.785] & 3.6$\times$ \\
1980--2023 (this revision) & 115 & \textbf{0.796 [0.753--0.837]} & \textbf{5.8$\times$} \\
\bottomrule
\end{tabular}
\end{center}

The first two rows have confidence intervals spanning chance; the audit's
diagnostic --- that sample size, not architecture, was the binding constraint
--- is thereby confirmed by intervention rather than argument. We still make
no operational skill claim: no comparison against SHIPS-RII or other
operational baselines has been performed (\S7).

\subsection{Spatial localization --- negative result, three independent angles}

The negative result on FuelMap localization survives the dataset expansion and
two additional controls:

\begin{enumerate}
  \item \textbf{Static TCHP validation} ($n = 226$ eligible test events,
        2022--2023 coverage): the FuelMap peak beats a random-point null
        ($p = 0.0003$) --- it tracks the storm --- but does \textbf{not}
        locate the TCHP peak better than the naive storm-centre baseline
        (median 539\,km vs 561\,km; closer in only 46\% of events; sign-flip
        permutation $p = 0.30$).
  \item \textbf{Dynamic displacement test:} the FuelMap's apparent collapse
        toward the storm centre during RI episodes, a candidate dynamic
        signal, was tested explicitly.
  \item \textbf{Physics-prior control:} repeating the displacement analysis
        on the pure enthalpy-flux prior --- no learned weights at all ---
        reproduces the same dynamic behavior. The ``signal'' is arithmetic of
        the physics prior, not learned skill.
\end{enumerate}

Na\"ive ``peak-of-field-in-window'' localization is further confounded by the
basin-scale ocean gradient (a majority of in-window TCHP maxima fall on the
window edge). \textbf{We therefore report FuelMap-based localization of the
energy source as unsupported}, and note that the model's demonstrable
classification skill (\S5.1) does not transfer to spatial attribution.

\subsection{ADT ocean input --- null result}

A controlled with/without-ADT ablation (same seed, data, epochs) shows ADT
does \textbf{not} improve RI classification --- it is marginally worse on
every split (full test ROC-AUC 0.599 $\rightarrow$ 0.578). With only 4--9 RI
positives per evaluation subset, the result is underpowered and inconclusive;
the honest reading is that ADT's value cannot be demonstrated in this data
regime. Separately, ADT does track TCHP at the storm centre (Spearman
$\rho \approx 0.30$, replicated for 2022 and 2023) --- but this merely
\textbf{reproduces known altimetry-OHC relationships}, it does not establish
new science. (This ablation predates the 1980--2023 expansion; a paired
feature ablation on the expanded dataset, with $\sim$800 positives of
statistical power, is planned as follow-up work.)

\subsection{Model-internal causal dependence --- positive but limited}

Occluding the FuelMap-identified region reduces predicted RI significantly
more than occluding a low-fuel control region (paired test,
$p \approx 10^{-6}$; $\sim$93--96\% of events). This shows the \emph{trained
model} relies on that region --- it does \textbf{not} show the region is the
true physical energy source. It is an interpretability check, not a physical
attribution.

\section{Discussion}

The audit clarifies what CycloneNet is. Its FuelMap is a
saliency/attention-like derived representation whose localization validity,
when tested against an independent ocean proxy, is not supported. Its
``physics guidance'' is weak supervision toward a heuristic prior, not a
physical law (and one of its physics terms is degenerate). Its ocean-input
extension reproduces established altimetry-OHC relationships without
improving skill. The framing of ``identifying the energy source feeding a
hurricane'' overstates the scope: the fuel is extracted locally at the
air--sea interface beneath the storm (not from a remote region), and mapping
it is already operational via TCHP/OHC and coupled models.

What remains genuinely useful is the artifact's transparency and the
discipline of testing each claim --- including reporting the negative ones.
We believe documenting these failure modes has value: physics-guided and
``interpretable'' cyclone-ML papers frequently assert that learned maps
recover physical quantities; here we show, with an explicit baseline, that a
plausible such map does not.

\section{Limitations}

\begin{itemize}
  \item \textbf{Surface-only input features.} The nine channels are
        surface/near-surface fields; the model sees neither the subsurface
        ocean reservoir that governs sustained RI nor the atmospheric column
        (e.g.\ deep-layer shear, mid-level humidity) that operational
        predictors exploit. Pressure-level extensions are implemented in the
        pipeline but not yet part of the released training set.
  \item \textbf{0.25$^{\circ}$ ERA5 resolution} cannot resolve inner-core
        structure; the model necessarily learns environmental favorability,
        not storm-scale dynamics.
  \item \textbf{No comparison against operational baselines (SHIPS-RII)} ---
        required before any operational skill claim; \S5.1's result is skill
        above chance, not skill above the state of practice.
  \item TCHP validation limited to 2022--2023 (gridded pre-2022 TCHP
        unavailable publicly), i.e.\ $n = 226$ of 2{,}679 test events.
  \item The heuristic prior is not a physical guarantee; the
        equation-consistency term is degenerate.
  \item The recall-first operating point yields precision 0.070 ---
        appropriate for the forensic screening framing, unusable as an alarm
        system.
\end{itemize}

\section{Correction record}

This manuscript corrects the prior version (``CycloneNet V2 \ldots{}
Atmospheric Singularity Mapping''). Specifically: (1) the released code did
not implement the physics-guided losses it described (weights defaulted to
0); this is fixed. (2) The threshold methodology in the code differed from
the text; this is reconciled. (3) The headline metrics (2{,}193 test samples,
ROC-AUC 0.83) were not reproducible from the public repository and derived
from a larger, unreleased dataset; this is \textbf{resolved in this revision}
by releasing the full 1980--2023 dataset and replacing the superseded numbers
with the reproducible results of \S5.1 --- the 0.83/0.905 figures should not
be cited. (4) The title and ``uniqueness/innovation'' framing overstated the
scope relative to existing literature and operational tools; the title is
revised and the framing corrected throughout. A detailed errata is provided
in the repository (\texttt{ERRATA.md}).

\section{Conclusion}

CycloneNet is an honest, reproducible, well-engineered hybrid physics+ML
pipeline for forensic RI analysis. It is \textbf{not} a novel method, a new
scientific finding, nor a tool that identifies the physical energy source of
hurricanes --- that capability already exists operationally. Its contribution
is engineering transparency plus a two-sided audit: on the positive side, RI
classification skill that is statistically distinguishable from chance and
fully reproducible from public artifacts (ROC-AUC 0.796, PR-AUC 5.8$\times$
chance, 115 test positives), obtained by fixing the audit-identified
bottleneck (sample size) rather than the architecture; on the negative side,
learned FuelMap localization does not beat a trivial baseline --- a result
now robust across a static TCHP comparison, a dynamic displacement test, and
a physics-prior control. We release it as a reproducible baseline and as a
record of what, in this approach, does and does not work.

\vspace{1em}
\noindent\emph{Code, data instructions, tests, and the errata are available
under CC BY-NC 4.0 at
\url{https://github.com/estefano-ferreira/cyclone-net}.}

\section*{Acknowledgements}

ERA5 (Copernicus C3S), IBTrACS (NOAA NCEI), TCHP (NOAA/AOML), and DUACS
sea-level data (Copernicus Marine). No external funding.

\end{document}
